% ===== UNIT 04: Contributions + paper organisation (1.3-1.4) =====
% Status: REWRITTEN  (change to REWRITTEN / REVIEWED when done; see CHECKLIST.md)
% Rules: style rewrite only. Keep every \cite, \ref, \label, number, equation and table value as is. Change the wording, not the content.
\subsection{Contributions}
The contributions of this work are as follows:
\begin{itemize}
  \item \textbf{C1: Reproducible ViT-B/16 pipeline and a
    recipe-versus-post-processing analysis.} We release a specific
    ViT-B/16 retrieval pipeline (CLIP initialisation, a tuned metric objective,
    an optimised augmentation pipeline, AQE $+$ $k$-reciprocal
    re-ranking~\cite{zhong2017rerank}, flip TTA) reaching \textbf{89.97\% mAP /
    97.79\% Rank-1} on VeRi-776. We test what drives this
    number: applying our identical post-processing to the official TransReID
    and CLIP-ReID checkpoints lifts them to \textbf{91.12\%} and
    \textbf{91.80\% mAP}, both slightly above ours, while all three
    are within 0.9\,pp without post-processing
    (82.11\,/\,82.99\,/\,82.22\%, Table~\ref{tab:matched}). Our recipe thus
    gives no advantage over these baselines at matched conditions: the ViT-B
    retrieval gains are because of post-processing, and we position
    Stream~1 as a reproducible baseline, not a superior recipe. Among
    single models reported with post-processing, MBR4B-LAI
    (92.1\% mAP~\cite{luo2023mbr4b}) is also higher.
  \item \textbf{C2: Marginal new best mAP via training-free fusion.}
    Score-level fusion with CLIP-SENet~\cite{lu2025clipsenet} achieves
    \textbf{93.32\% mAP / 98.45\% Rank-1}, $+0.42$\,pp mAP above the best
    previously published result (92.9\%~\cite{lu2025clipsenet}) and
    $+1.88$\,pp above our own reproduction of that stream (91.44\%). Rank-1
    (98.45\%) is competitive with but does not exceed CLIP-SENet standalone
    Rank-1 (98.7\%). Because the $+0.42$\,pp margin is small and our fusion reuses that model's
    released checkpoint, we present it as a demonstration of feature
    complementarity, not a decisive state of the art. A paired bootstrap test
    shows that the gain over the best single stream is significant
    (Section~\ref{sec:holdout}).
  \item \textbf{C3: Same-setup ablation path.} An in depth
    ablation traces the complete performance trajectory from the 82.3\%
    TransReID baseline to 93.3\% mAP under a single consistent post-processing
    protocol, separating post-processing contributions from recipe
    contributions instead of conflating them.
  \item \textbf{C4: Reproducible public release with validation and cost
    analysis.} All checkpoints (SHA-256 pinned), evaluation notebooks,
    configuration files, per-query evaluation dumps, and the post-processing
    pipeline are publicly released. Every reported number is the deterministic
    re-evaluation of a frozen, hash-pinned checkpoint and is thus
    exactly re-computable from the release. From-scratch retraining is
    reproducible within a $\pm0.55$\,pp seed band
    (Section~\ref{sec:seed_variance}). Headline numbers use hyperparameters
    selected on the standard evaluation split (disclosed, not implicit), with
    held-out validation bounding the resulting optimism to $\leq$0.2\,pp mAP
    (Section~\ref{sec:holdout}), and the latency and memory of every pipeline
    stage (extraction, TTA, AQE, re-ranking, fusion) are measured for
    deployment planning (Section~\ref{sec:cost}).
  \item \textbf{C5: Cross-dataset transfer study.} A zero-shot
    VeRi-776$\rightarrow$VeRi-Wild evaluation across all three official test
    protocols (with paired-bootstrap CIs and released per-query dumps) shows
    the ViT-B stream transfers competitively with published cross-domain
    methods that use over twenty times more source identities or adapt using
    unlabelled target data, while the two-stream fusion does not
    transfer at any scale: evidence, supported by a CityFlowV2 MTMC port,
    that its VeRi-776 advantage is domain-specific
    (Section~\ref{sec:veriwild}).
\end{itemize}

\subsection{Paper organisation}
Section~\ref{sec:related} reviews related work across CNN-era methods,
transformer-based methods, and post-processing techniques.
Section~\ref{sec:method} describes the full methodology including both streams
and the fusion strategy. Section~\ref{sec:setup} details the experimental
setup. Section~\ref{sec:results} presents quantitative results and ablations.
Section~\ref{sec:discussion} discusses findings and limitations.
Section~\ref{sec:conclusion} concludes.
